% Chapter 10: the PC program "winder".
\chapter{The PC program (winder)}\label{ch:pc}
\index{winder program}\index{PC program}

\textbf{winder} is a small program for a Windows, Mac, Linux or Chromebook
computer. It shows the machine's state, calibrates it and edits the
settings. It runs a page in your web browser; nothing is sent over the
internet. You do \emph{not} need it to wind paper.

\section{What you need}\label{sec:adapter}
\index{serial adapter}\index{USB serial adapter}
\begin{itemize}
  \item A computer.
  \item A \textbf{USB serial adapter} with 5 volt logic (common chips: CH340,
        CP210x, FTDI). Most computers have the driver built in.
  \item Three wires from the adapter to the controller board:
\end{itemize}
\begin{center}
\begin{tabular}{@{}lll@{}}
\toprule
\textbf{Adapter pin} & \textbf{Controller pin} & \\
\midrule
TX  & D0 (RX) & \\
RX  & D1 (TX) & \\
GND & GND & \\
\bottomrule
\end{tabular}
\end{center}
\begin{warning}
Do not connect the adapter's 5 V or 3.3 V power pin. The machine has its own
power. The controller's own USB port is \emph{not} used by the winder program
(it is only for loading firmware, Chapter~\ref{ch:flashing}).
\end{warning}

\section{Download}
Get the program from the release page
\url{https://github.com/lognd/stpone/releases/tag/winder-latest}. Pick the file
for your computer:
\begin{center}
\renewcommand{\arraystretch}{1.3}
\begin{tabularx}{\textwidth}{@{}>{\raggedright\arraybackslash}p{0.30\textwidth}>{\raggedright\arraybackslash}X@{}}
\toprule
\textbf{Computer} & \textbf{File} \\
\midrule
Windows & \texttt{winder-windows-x86\_64.exe} \\
Mac with Intel chip & \texttt{winder-macos-x86\_64.tar.gz} \\
Mac with Apple chip (M1 and later) & \texttt{winder-macos-aarch64.tar.gz} \\
Linux PC & \texttt{winder-linux-x86\_64.tar.gz} \\
Linux on ARM, Raspberry Pi & \texttt{winder-linux-aarch64.tar.gz} \\
Chromebook & see Section~\ref{sec:chromebook} \\
\bottomrule
\end{tabularx}
\end{center}
The release is rebuilt automatically whenever the software passes all its
tests, so ``latest'' is always the current version.

\section{Start the program}
\begin{description}[style=nextline,leftmargin=1.5em]
  \item[Windows] Double-click \texttt{winder-windows-x86\_64.exe}. If Windows
    warns about an unknown publisher, choose to run it anyway.
  \item[Mac and Linux] Open a terminal in the download folder and type (for
    example)\par
    \texttt{tar -xzf winder-linux-x86\_64.tar.gz}\par
    \texttt{./winder}
\end{description}
The program opens your web browser at \texttt{http://127.0.0.1:8765}. If that
port is busy it picks another and prints the address. If it cannot open a
browser it writes the address into a file \texttt{winder-open-me.txt}; open
that address yourself. To quit, close the program's window (or press
Ctrl+C in the terminal).

\section{Chromebook}\label{sec:chromebook}
\index{Chromebook}
\begin{enumerate}
  \item Open \ui{Settings > Advanced > Developers}, find \ui{Linux development
        environment} and click \ui{Turn on}.
  \item In the same place, open \ui{Manage USB devices} and switch the serial
        adapter on for Linux.
  \item In the Linux terminal, type \texttt{uname -m}. If it says
        \texttt{x86\_64} use the Linux x86\_64 file, if it says
        \texttt{aarch64} use the aarch64 file.
  \item Unpack and start it as for Linux above. You may need to add your user
        to the \texttt{dialout} group: \texttt{sudo usermod -aG dialout \$USER},
        then restart Linux.
  \item Open \texttt{localhost:8765} in the Chrome browser.
\end{enumerate}
\begin{note}
The Chromebook steps have not yet been tested on a real Chromebook.
\end{note}

\section{Connection check (doctor)}\label{sec:doctor}
\index{connection check}\index{doctor}
If anything does not work, open the \ui{Advanced} tab and look at
\ui{Connection check}, or run \texttt{./winder doctor} in a terminal. It tests,
in order:
\begin{enumerate}
  \item \ui{USB serial adapter found}
  \item \ui{Allowed to use the adapter}
  \item \ui{Winder answers}
  \item \ui{Firmware is compatible}
  \item \ui{Saved settings are valid}
  \item \ui{A web browser can be opened}
\end{enumerate}
Each failed line says what to do. Some have a \ui{Fix it} button. Click
\ui{Check again} after fixing something.

\section{The screens of the program}
\index{winder program!tabs}
A \ui{Stop motor} button is always at the top of the page. It stops anything
the program started (calibration, Test mode runs). It cannot stop a wind
started from the panel; use \btn{1} or \btn{2} for that.

\begin{description}[style=nextline,leftmargin=1.5em]
  \item[\ui{Home}] Machine status, whether the motor is on, the length wound,
    and the \ui{Calibrate the machine} button.
  \item[\ui{Calibrate}] The step-by-step calibration (Chapter~\ref{ch:calibration}).
  \item[\ui{Settings}] The five editable settings (Chapter~\ref{ch:settings}),
    with \ui{Save to the machine}, \ui{Undo all} and
    \ui{Reset to factory settings}.
  \item[\ui{Advanced}] For technicians: the \ui{Test mode} switch,
    \ui{Move the machine by hand} (\ui{Motor on/off}, \ui{Set length wound to
    zero}, \ui{Run} a set length), \ui{Last test winds} (\ui{Start test winds},
    \ui{Clear table}, \ui{Apply suggestions}), \ui{Machine state},
    \ui{Connection check} and \ui{What happened} (a log of recent events).
\end{description}

\section{Test mode}\label{sec:testmode}
\index{test mode}
Calibration and the \ui{Advanced} controls put the machine in \textbf{Test
mode}. The screen shows \shows{DEBUG: HOST CONTROL} and \shows{panel buttons
locked}; \led{2} and \led{4} light. In Test mode:
\begin{itemize}
  \item The panel cannot start a wind.
  \item \btn{1} or \btn{2} on the panel stops the motor and ends Test mode.
  \item If the PC stops talking to the machine for 5 seconds (cable pulled,
        program crashed), the machine stops the motor and leaves Test mode by
        itself.
  \item If you close the browser page, the program stops the motor and ends
        Test mode after about 12 seconds.
\end{itemize}

\section{Command line options}
For technicians. Type \texttt{winder --help} for the full list.
\begin{center}
\small
\begin{tabularx}{\textwidth}{@{}>{\ttfamily}l>{\raggedright\arraybackslash}X@{}}
\toprule
\textrm{\textbf{Option}} & \textbf{Meaning} \\
\midrule
--port & Use this serial port instead of searching. \\
--baud & Serial speed (the machine uses 115200). \\
--http-port & Web page port (default 8765). \\
--no-browser & Do not open a browser. \\
--list-ports & List serial ports and exit. \\
--fake & Pretend machine, for practice without hardware. \\
--true-coast & With --fake: the coast, in counts, the pretend machine has. \\
doctor & Run the connection check in the terminal. \\
calibrate & Calibrate from the terminal (\texttt{--lengths}, \texttt{--measured}, \texttt{--save}). \\
\bottomrule
\end{tabularx}
\end{center}
